\chapter{Markov Decision Process (MDP)}

\lecture{Jan. 28}

\section{Q-value function}
Value function of policy $\pi$ is defined as
\[
    V_\pi(s) = \E[\sum_{t=0}^{\infty} \gamma^t r(s_t, a_t)|s_0 = s]
\]
Action-value function for policy $\pi$ is defined as
\[
    Q_\pi(s, a) = \E [\sum_{t=0}^{\infty} \gamma^t r(s_t, a_t)|s_0=s, a_0=a]
\]
It is easy to show
\[
    Q_\pi(s, a) = r(s, a) + \gamma \sum_{s'} P(s' | s, a) V_\pi (s')
\]
this is the instantaneous reward for $(s,a)$ and the expected value of the next state.

\section{Value Policy Iteration}
\subsubsection{Fixed point}
%\def 
Let $f : \mathcal{U} \to \mathcal{U}$. $u^*$ is a fixed point of $f(\cdot)$ if $f(u^*) = u^*$.

%\def 
Let $f : \mathcal{U} \to \mathcal{U}$ and let $\norm{\cdot}$ be a norm defined on $\mathcal{U}$. We say $f(\cdot)$ is a contraction mapping (or simply,
contraction) if $\norm{f(u) - f(v)} \leq \alpha\norm{u-v}$ for an $\alpha \in (0, 1)$ and for all $u, v \in \mathcal{U}$.

%\thm
Banach Fixed Point Theorem.
\[
    \norm{u_{n+1} - u_n} \leq \alpha^n \norm{u_1 - u_0}
\]
converges when $\alpha \leq 1$. We can show that the sequence $\{u_0, u_1, u_2, \dots \}$ is a Cauchy sequence. Since $\mathcal{U}$ is closed and bounded,
the sequence converges, i.e., there exists a limit $u^* = \lim_{k \to \infty} u_k$. Since $f(\cdot)$ is a contraction,
it is also a continuous function. Then, $u^* = \lim_{k\to\infty} u_k = f(\lim_{k\to\infty} u_k) = f(u^*)$.

Proof of uniqueness: suppose there exists another fixed point $\bar{u}^* \neq u^*$. Then,
\[
    \norm{u^* - \bar{u}^*} = \norm{f(u^*) - f(\bar{u}^*)} \leq \alpha \norm{u^* - \bar{u}^*}.
\]
This is a contradiction, so $u^* = \bar{u}^*$.

\section{Bellman Optimality Equation}
Let $V^*$ be the optimal value function. Then, $V^*$ satisfies the equation
\[
    V^*(s) = \max_{a} \left( r(s,a) + \gamma \sum_{s'} P(s' | s,a) V^*(s') \right)
\]
Suppose the current state is $s$ and we take action $a$. From the next state onwards, we follow the optimal policy $\pi^*$. What
is the 'value' of state $s$ (expected cumulative discounted reward) under this procedure?
\[
    V(s) = r(s,a) + \gamma \E_{s' \sim P(\cdot | s,a)} [V^*(s')]
\]
reward for $(s,a)$ plus $gamma$ times the expected value of the next state under the optimal policy $\pi^*$. The best action we can take
in the current state $s$ is the action which maximizes the value $V(s)$. The 'value' of state $s$ if we take this action is
\[
    \max_{a} (r(s,a) + \gamma \E_{s\sim P(\cdot | s,a)} [V^*(s')])
\]
\section{Value Iteration}
Bellman Operator $T: \R^{|\mathcal{S}|} \to \R^{|\mathcal{S}|}$ is defined as
\[
    (TV)(s) = \max_{a} (r(s, a) + \gamma \E_{s' \sim P(\cdot | s, a)}[V(s')])
\]
Aim to design feedback, with $V_{k+1} = TV_k$, for which we want to prove converge.

Proposition: Bellman operator $T$ is a contraction with respect to $\norm{\cdot}_\infty$, i.e., for any $V_1, V_2 \in \mathcal{R}^{|S|}$,
$\norm{TV_1 - TV_2}_{\infty} \leq \gamma \norm{V_1 - V_2}_{\infty}$.

%\lemma
A useful result for the proof. Let $g,h$ be any two real-valued function. Then,
\[
    | \max_{a} g(x,a) - \max_{a} h(x,a) | \leq \max_{a} |g(x, a) - h(x, a)|
\]
Proof (of contraction property)
\begin{align*}
    |(TV_1)(s) - (TV_2)(s)| & = |\max_{a \in \mathcal{A}}(r(s,a) + \gamma \E_{s' \sim P(\cdot | s, a)}[V_1(s')]) - \max_{a \in \mathcal{A}}(r(s,a) + \gamma \E_{s' \sim P(\cdot | s, a)}[V_2(s')])| \\
                            & \leq \max_{a \in \mathcal{A}} |(r(s,a) + \gamma \E_{s' \sim P(\cdot | s, a)}[V_1(s')]) - (r(s,a) + \gamma \E_{s' \sim P(\cdot | s, a)}[V_2(s')])|                     \\
                            & = \max_{a \in \mathcal{A}} |\gamma \E_{s' \sim P(\cdot | s, a)}[V_1(s')] - \gamma \E_{s' \sim P(\cdot | s, a)}[V_2(s')])|                                             \\
                            & \leq \max_{a \in \mathcal{A}} \gamma \E_{s' \sim P(\cdot | s,a)} [|V_1(s') - V_2(s')|]                                                                                \\
                            & \leq \max_{a \in \mathcal{A}} \gamma \E_{s' \sim P(\cdot| s,a)} [\norm{V_1 - V_2}_{\infty}] = \gamma \norm{V_1 - V_2}_{\infty}
\end{align*}
this implies
\[
    \norm{TV_1 - TV_2}_{\infty} \leq \gamma \norm{V_1 - V_2}_{\infty}
\]